\chapter{Rational numbers}
\label{ch:rationals}

In $\Z$, $2 \cdot x = 1$ has no solution. We repeat the trick of
\cref{ch:integers}, now for division: a rational number is what a quotient of
integers \emph{would} be. As a stand-in we use the pair $(a, b)$ with $b \neq 0$.
Two pairs mean the same quotient when \enquote{$a$ over $b$} equals \enquote{$c$
over $d$}, that is, when $a \cdot d = c \cdot b$. That equation only uses
multiplication, which we have.

\section{Construction}

\begin{definition}
\label{def:rational-relation}
Let $\Z^\ast \coloneqq \Z \setminus \set{0}$. On $\Z \times \Z^\ast$ let
\[
  (a, b) \approx (c, d) \defiff a \cdot d = c \cdot b.
\]
\register{rationals}{$\Z^\ast$}{the nonzero integers}
\end{definition}

\begin{lemma}
\label{lem:rational-equivalence}
$\approx$ is an equivalence relation.
\end{lemma}

\begin{proof}
Reflexive and symmetric are clear. Transitive: let $a \cdot d = c \cdot b$ and $c
\cdot f = e \cdot d$. Multiplying the first by $f$ gives $a \cdot d \cdot f = c
\cdot b \cdot f = (c \cdot f) \cdot b = e \cdot d \cdot b$, so $(a \cdot f) \cdot d
= (e \cdot b) \cdot d$. Since $d \neq 0$, cancellation in $\Z$
(\cref{thm:integral-domain}(2)) gives $a \cdot f = e \cdot b$.
\end{proof}

This is where \cref{ch:integers} pays off: without \enquote{no zero divisors},
$\approx$ would not be transitive.

\begin{definition}[Rational numbers]
\label{def:rationals}
$\Q \coloneqq (\Z \times \Z^\ast)/{\approx}$. Its elements are the \emph{rational
numbers}. The class of $(a, b)$ is $[(a, b)]$, with \emph{numerator} $a$ and
\emph{denominator} $b$.
\register{rationals}{$\Q$}{the set of rational numbers}
\register{rationals}{$[(a, b)]$}{the rational number \enquote{$a$ over $b$};
replaced by $a / b$ after \cref{thm:rationals-embedding}}
\end{definition}

So $[(a, b)] = [(c, d)] \Leftrightarrow a \cdot d = c \cdot b$.

\begin{lemma}[Expanding]
\label{lem:expand}
$[(a \cdot k, b \cdot k)] = [(a, b)]$ for $k \neq 0$. In particular $[(-a, -b)] =
[(a, b)]$, so every rational number has a representative whose denominator lies in
$\N$.
\end{lemma}

\begin{proof}
$(a \cdot k) \cdot b = a \cdot (b \cdot k)$. Take $k = -1$ for the second claim,
and use \cref{cor:integer-normal-form} for the third.
\end{proof}

\section{Arithmetic}

\begin{definition}[Operations on $\Q$]
\label{def:rational-operations}
\begin{align*}
  [(a, b)] + [(c, d)] &\coloneqq [(a \cdot d + c \cdot b,\; b \cdot d)], \\
  -[(a, b)] &\coloneqq [(-a, b)], \\
  [(a, b)] \cdot [(c, d)] &\coloneqq [(a \cdot c,\; b \cdot d)],
\end{align*}
with $0 \coloneqq [(0, 1)]$ and $1 \coloneqq [(1, 1)]$. The denominators are
allowed: $b, d \neq 0$ gives $b \cdot d \neq 0$ by \cref{thm:integral-domain}(1).
As in $\Z$, the symbols are shared with the smaller number systems
(\cref{rule:identification,thm:rationals-embedding}).
\register{rationals}{$x + y$, $-x$, $x \cdot y$, $0$, $1$}{the ring operations of
$\Q$ (extend those of $\Z$)}
\end{definition}

\begin{lemma}
\label{lem:rational-well-defined}
The operations of \cref{def:rational-operations} are well-defined.
\end{lemma}

\begin{proof}
Let $a \cdot b' = a' \cdot b$. We change the first argument; both formulas are
symmetric in the two arguments, so that is enough.

\emph{Addition.} We need $(a \cdot d + c \cdot b) \cdot (b' \cdot d) = (a' \cdot d
+ c \cdot b') \cdot (b \cdot d)$. Expanded, the left side is $a \cdot b' \cdot d
\cdot d + c \cdot b \cdot b' \cdot d$ and the right side is $a' \cdot b \cdot d
\cdot d + c \cdot b' \cdot b \cdot d$. They are equal because $a \cdot b' = a'
\cdot b$.

\emph{Multiplication.} $(a \cdot c) \cdot (b' \cdot d) = (a \cdot b') \cdot c
\cdot d = (a' \cdot b) \cdot c \cdot d = (a' \cdot c) \cdot (b \cdot d)$.

\emph{Negation.} $(-a) \cdot b' = -(a \cdot b') = -(a' \cdot b) = (-a') \cdot b$.
\end{proof}

\begin{definition}[Field]
\label{def:field}
A \emph{field} is a commutative ring with $0 \neq 1$ in which every $x \neq 0$ has
a \emph{multiplicative inverse} $x^{-1}$ with $x \cdot x^{-1} = 1$. For $y \neq 0$,
the \emph{quotient} is $x / y \coloneqq x \cdot y^{-1}$.
\register{rings}{$x^{-1}$}{multiplicative inverse of $x \neq 0$ in a field}
\register{rings}{$x / y$}{$x \cdot y^{-1}$}
\end{definition}

\begin{theorem}
\label{thm:rationals-field}
$\Q$ is a field. The inverse of $[(a, b)] \neq 0$ is $[(b, a)]$.
\end{theorem}

\begin{proof}
Let $x = [(a, b)]$, $y = [(c, d)]$ and $z = [(e, f)]$. Commutativity of both
operations is built into the formulas.

\emph{Associativity.} $(x + y) + z$ and $x + (y + z)$ both equal $[(a \cdot d \cdot
f + c \cdot b \cdot f + e \cdot b \cdot d,\; b \cdot d \cdot f)]$, and $(x \cdot y)
\cdot z$ and $x \cdot (y \cdot z)$ both equal $[(a \cdot c \cdot e,\; b \cdot d
\cdot f)]$.

\emph{Neutral elements.} $x + 0 = [(a \cdot 1 + 0 \cdot b,\; b \cdot 1)] = x$, and
$x \cdot 1 = [(a \cdot 1,\; b \cdot 1)] = x$.

\emph{Additive inverse.} $x + (-x) = [(a \cdot b + (-a) \cdot b,\; b \cdot b)] =
[(0,\; b \cdot b)] = [(0, 1)]$, since $0 \cdot 1 = 0 \cdot (b \cdot b)$.

\emph{Distributivity.} Using \cref{lem:expand} with $k = b$:
\[
  x \cdot y + x \cdot z = [(a \cdot c \cdot b \cdot f + a \cdot e \cdot b \cdot
  d,\; b \cdot d \cdot b \cdot f)] = [(a \cdot (c \cdot f + e \cdot d),\; b \cdot d
  \cdot f)] = x \cdot (y + z).
\]

\emph{$0 \neq 1$.} $[(0, 1)] = [(1, 1)]$ would mean $0 = 1$ in $\Z$.

\emph{Multiplicative inverse.} First, $[(a, b)] = 0$ if and only if $a \cdot 1 = 0
\cdot b$, if and only if $a = 0$. So if $x \neq 0$, then $a \neq 0$, and $[(b, a)]$
is a rational number with $[(a, b)] \cdot [(b, a)] = [(a \cdot b,\; b \cdot a)] =
[(1, 1)]$.
\end{proof}

All of \cref{lem:ring-laws} holds in $\Q$, since only ring laws were used there.
A field has no zero divisors: if $x \cdot y = 0$ and $x \neq 0$, then $y = x^{-1}
\cdot x \cdot y = 0$.

\section{\texorpdfstring{$\Z$}{Z} inside \texorpdfstring{$\Q$}{Q}}

\begin{theorem}[Embedding]
\label{thm:rationals-embedding}
Let $j \colon \Z \to \Q$ be the map with $a \via{j} [(a, 1)]$. Then $j$ is
injective, $0 \via{j} 0$, $1 \via{j} 1$, and for all $a, c \in \Z$, $b \in
\Z^\ast$ and $x, y, u \in \Q$:
\[
  a \via{j} x \land c \via{j} y \;\Rightarrow\; a + c \via{j} x + y \;\land\; a
  \cdot c \via{j} x \cdot y,
  \qquad
  a \via{j} x \land b \via{j} u \;\Rightarrow\; [(a, b)] = x / u.
\]
\end{theorem}

\begin{proof}
$[(a, 1)] = [(a', 1)]$ means $a = a'$. Next, $[(a, 1)] + [(c, 1)] = [(a + c, 1)]$
and $[(a, 1)] \cdot [(c, 1)] = [(a \cdot c, 1)]$. Finally the inverse of $[(b,
1)]$ is $[(1, b)]$ by \cref{thm:rationals-field}, and $[(a, 1)] \cdot [(1, b)] =
[(a, b)]$.
\end{proof}

As before, \emph{we identify every $a \in \Z$ with $[(a, 1)]$}, so $\N \subseteq \Z
\subseteq \Q$ (\cref{rule:identification}). The theorem then says $[(a, b)] = a /
b$: the class is the quotient. By \cref{rule:abbreviation}, from now on rationals
are written as quotients and the class notation is not used again. In this
notation, $a / b = c / d \Leftrightarrow a \cdot d = c \cdot b$.

\section{Order}

\begin{definition}[Order on $\Q$]
\label{def:rational-order}
A rational number is \emph{nonnegative} if it equals $a / b$ for some $a, b \in \N$
with $b \neq 0$. For $x, y \in \Q$ let
\[
  x \le y \defiff y - x \text{ is nonnegative}.
\]
\register{rationals}{$x \le y$, $x < y$}{order of rationals: $y - x = a / b$ with
$a, b \in \N$ (extends the order of $\Z$)}
\end{definition}

\begin{theorem}[$\Q$ is an ordered field]
\label{thm:rationals-ordered}
$\le$ is a total order on $\Q$; $x \le y \Rightarrow x + z \le y + z$; and $0 \le
x \land 0 \le y \Rightarrow 0 \le x \cdot y$. On $\Z$ it agrees with the order of
\cref{ch:integers}.
\end{theorem}

\begin{proof}
$0 \le x$ means that $x$ is nonnegative. The nonnegative rationals have three
properties:
\begin{itemize}
  \item \emph{They are closed under $+$ and $\cdot$:} for $a, b, c, d \in \N$ with
    $b, d \neq 0$, both $a / b + c / d = (a \cdot d + c \cdot b) / (b \cdot d)$ and
    $(a / b) \cdot (c / d) = (a \cdot c) / (b \cdot d)$ are quotients of natural
    numbers.
  \item \emph{For every $x$, $x$ or $-x$ is nonnegative:} write $x = a / b$ with $b
    \in \N$, $b \neq 0$ (\cref{lem:expand}). If $a \in \N$, $x$ is nonnegative;
    otherwise $-a \in \N$ and $-x = (-a) / b$ is.
  \item \emph{Only $0$ is both:} if $x = a / b$ and $-x = c / d$ with $a, b, c, d
    \in \N$, then $a / b = (-c) / d$, so $a \cdot d = -(c \cdot b)$, so $a \cdot d
    + c \cdot b = 0$ in $\N$. Then $a \cdot d = 0$ (\cref{thm:addition-laws}(6)),
    so $a = 0$ and $x = 0$.
\end{itemize}
The proof of \cref{thm:integers-ordered}(1)--(3) used only these three properties
of $\N \subseteq \Z$, so it goes through word for word.

On $\Z$: if $a \in \N$, then $a = a / 1$ is nonnegative. Conversely, if $a \in \Z$
and $a = c / d$ with $c, d \in \N$, $d \neq 0$, then $a \cdot d = c \in \N$. If $a$
were negative, $-a \in \N \setminus \set{0}$ and $(-a) \cdot d = -c$ would be a
nonzero natural number with $c + (-a) \cdot d = 0$, which is impossible
(\cref{thm:addition-laws}(6)). So $a \in \N$.
\end{proof}

\section{Two properties of \texorpdfstring{$\Q$}{Q}}

\begin{theorem}[Archimedean property]
\label{thm:archimedean}
For every $x \in \Q$ there is an $n \in \N$ with $x < n$.
\end{theorem}

\begin{proof}
Write $x = a / b$ with $b \in \N$, $b \neq 0$, so $1 \le b$. If $a \le 0$, then $x
\le 0 < 1$. If $0 < a$, take $n \coloneqq a + 1$. Then
\[
  n - x = \bigl( (a + 1) \cdot b - a \bigr) / b = \bigl( a \cdot (b - 1) + b \bigr)
  / b,
\]
a quotient of natural numbers whose numerator is not $0$. So $x < n$.
\end{proof}

\begin{theorem}[Density]
\label{thm:density}
If $x < y$ in $\Q$, then $x < (x + y) / 2$ and $(x + y) / 2 < y$. So between any two
rational numbers there is another one, and hence infinitely many.
\end{theorem}

\begin{proof}
$0 < 2$, so $0 < 2^{-1}$. From $x < y$ we get $2 \cdot x = x + x < x + y$ and $x + y
< y + y = 2 \cdot y$. Multiply by $2^{-1}$.
\end{proof}

\section{The gap}

$\Q$ looks complete: it is a field, it is ordered, and it is dense. But it has
holes.

\begin{theorem}[The square root of $2$ is irrational]
\label{thm:sqrt2}
There is no $q \in \Q$ with $q \cdot q = 2$.
\end{theorem}

\begin{proof}
Suppose $q \cdot q = 2$. Since $(-q) \cdot (-q) = q \cdot q$, we may assume $0 <
q$. Then $1 < q$ and $q < 2$: $q \le 1$ would give $q \cdot q \le q \le 1$, and $2
\le q$ would give $4 \le q \cdot q$.

Write $q = a / b$ with $a, b \in \N$, $b \neq 0$. Then $b \cdot q = a$ is a natural
number, so the set
\[
  X \coloneqq \setof{m \in \N}{m \neq 0 \land m \cdot q \in \N}
\]
is nonempty. By \cref{thm:well-ordering} it has a least element $m$. Let $m'
\coloneqq m \cdot q - m = m \cdot (q - 1)$. Then:
\begin{itemize}
  \item $m'$ is an integer, as a difference of two natural numbers, and $0 < m'$
    because $1 < q$. So $m' \in \N$ and $m' \neq 0$.
  \item $m' \cdot q = m \cdot q \cdot q - m \cdot q = 2 \cdot m - m \cdot q$ is an
    integer, and $0 < m' \cdot q$ because $0 < m'$ and $0 < q$. So $m' \cdot q \in
    \N$.
  \item $m' < m$, because $q - 1 < 1$.
\end{itemize}
So $m' \in X$ and $m' < m$, contradicting the choice of $m$.
\end{proof}

The proof uses only order and well-ordering; it needs no divisibility. So the set
$\setof{q \in \Q}{0 < q \land q \cdot q < 2}$ is bounded above (by $2$), but its
least upper bound would have to be a square root of $2$, and there is none in $\Q$.

This is where the real numbers come in. They are constructed by filling exactly
these gaps, using sequences of rational numbers, i.e.\ maps $\N \to \Q$, or cuts
of $\Q$. From there come limits, series, functions and calculus.

\section{Summary of the number constructions}

\begin{center}
\begin{tabular}{@{}lllll@{}}
  \toprule
  Set & Built from & Element & Becomes & Gains \\
  \midrule
  $\N$ & $\varnothing$ and successors & $\set{0, 1, \dots}$ & $n$ & counting,
    induction, recursion \\
  $\Z$ & classes of pairs in $\N \times \N$ & $[(a, b)]$ & $a - b$ & subtraction \\
  $\Q$ & classes of pairs in $\Z \times \Z^\ast$ & $[(a, b)]$ & $a / b$ & division \\
  $\R$ & still to come & & & limits, completeness \\
  \bottomrule
\end{tabular}
\end{center}

\begin{exercises}
  \item Show that $\Q$ is the smallest field containing $\Z$, in this sense: if $F$
    is a field and $\varphi \colon \Z \to F$ is injective and preserves $+$, $\cdot$
    and $1$, then there is exactly one injective map $\psi \colon \Q \to F$ that
    preserves $+$ and $\cdot$ and agrees with $\varphi$ on $\Z$.
  \item Show that for $x, y \in \Q$ with $0 < x$ there is $n \in \N$ with $y < n
    \cdot x$.
  \item Adapt the proof of \cref{thm:sqrt2} to show there is no rational $q$ with
    $q \cdot q = 3$.
  \item Let $S \coloneqq \setof{q \in \Q}{0 < q \land q \cdot q < 2}$. Show that for
    every $s \in S$ there is a larger element of $S$. (Try $s + h$ with a suitable
    small rational $0 < h$.)
\end{exercises}
